% =============================================================================
%  Atto 1 — Il rover come case study  (~3 min, 1 slide)
% =============================================================================

\section{The platform under inspection}

\begin{frame}{The MiviaRover, in one slide}
  \begin{columns}[T,onlytextwidth]
    \column{0.46\textwidth}
      \textbf{Custom outdoor UGV} built on a Lynxmotion \emph{Rugged~Rover}
      chassis, configured as a \textbf{tank-like differential drive}.

      \vspace{0.6em}
      \textbf{Numbers that matter for this talk}
      \begin{itemize}\setlength\itemsep{2pt}
        \item \textbf{4 brushed DC motors}, one per wheel
              (up to \(\sim\!160\)~RPM)
        \item \textbf{48-CPR quadrature encoder} on each shaft\\
              \footnotesize\textcolor{black!60}{(24-CPR $\times$ 2 channels;
              gear ratio 51:1)}
        \item \textbf{Custom PCB} — H-bridges, power, sensor I/O
        \item \textbf{Low-level} brain: \textbf{STM32F767ZI} + FPU,
              FreeRTOS
        \item \textbf{High-level} brain: \textbf{Jetson AGX Orin}
        \item Exteroception: \textbf{ZED~2i}; absolute fix:
              \textbf{Qorvo MDEK1001 UWB}
      \end{itemize}

      \vspace{0.4em}
      \footnotesize\textcolor{black!60}{Two halves talking over a
      \textbf{CAN} bus. The split is the whole story of this talk.}

    \column{0.54\textwidth}
      \centering
      \includegraphics[width=0.98\linewidth]{MiviaRoverPhysicalView.png}\\
      \scriptsize \emph{Hard-real-time hardware in blue, soft-real-time in
      orange, CAN bus in between. External sensors over USB.}
  \end{columns}
\end{frame}
